\documentclass[11pt,a4paper]{article}
% Wspólny szablon dla kart supermocy (jedna karta = jedna strona A4)
% Użycie w karcie: \documentclass[11pt,a4paper]{article} % Wspólny szablon dla kart supermocy (jedna karta = jedna strona A4)
% Użycie w karcie: \documentclass[11pt,a4paper]{article} % Wspólny szablon dla kart supermocy (jedna karta = jedna strona A4)
% Użycie w karcie: \documentclass[11pt,a4paper]{article} \input{../preambula} \begin{document}

% Silnik tectonic (XeTeX): polskie znaki przez fontspec, nie przez inputenc/T1
\usepackage{fontspec}
\setmainfont{Arial}
\usepackage[polish]{babel}
\usepackage[a4paper,margin=1.5cm]{geometry}
\usepackage{amsmath,amssymb}
\usepackage{tikz}
\usetikzlibrary{arrows.meta,positioning,calc,shapes.geometric}
\usepackage[most]{tcolorbox}
\usepackage{enumitem}
\usepackage{booktabs}
\usepackage{xcolor}

\pagestyle{empty}
\setlength{\parindent}{0pt}
\setlength{\parskip}{0.4em}

\definecolor{cel}{HTML}{1F4E79}
\definecolor{regula}{HTML}{E8F1FA}
\definecolor{pulapka}{HTML}{FDECEA}
\definecolor{okres}{HTML}{C0392B}
\definecolor{skonczony}{HTML}{1E8449}

% Ramka na regułę
\newtcolorbox{regula}[1][]{
  colback=regula, colframe=cel, boxrule=0.8pt, arc=2pt,
  left=6pt, right=6pt, top=4pt, bottom=4pt, title=#1,
  fonttitle=\bfseries, coltitle=white
}

% Ramka na pułapkę z kartkówki
\newtcolorbox{pulapka}[1][]{
  colback=pulapka, colframe=okres, boxrule=0.8pt, arc=2pt,
  left=6pt, right=6pt, top=4pt, bottom=4pt, title=#1,
  fonttitle=\bfseries, coltitle=white
}

% Nagłówek karty: numer, tytuł, cel
\newcommand{\naglowekKarty}[3]{%
  {\LARGE\bfseries\color{cel} Karta #1: #2}\\[2pt]
  {\large Cel: #3}\par
  \vspace{0.3em}\noindent\rule{\linewidth}{1pt}\par
}
 \begin{document}

% Silnik tectonic (XeTeX): polskie znaki przez fontspec, nie przez inputenc/T1
\usepackage{fontspec}
\setmainfont{Arial}
\usepackage[polish]{babel}
\usepackage[a4paper,margin=1.5cm]{geometry}
\usepackage{amsmath,amssymb}
\usepackage{tikz}
\usetikzlibrary{arrows.meta,positioning,calc,shapes.geometric}
\usepackage[most]{tcolorbox}
\usepackage{enumitem}
\usepackage{booktabs}
\usepackage{xcolor}

\pagestyle{empty}
\setlength{\parindent}{0pt}
\setlength{\parskip}{0.4em}

\definecolor{cel}{HTML}{1F4E79}
\definecolor{regula}{HTML}{E8F1FA}
\definecolor{pulapka}{HTML}{FDECEA}
\definecolor{okres}{HTML}{C0392B}
\definecolor{skonczony}{HTML}{1E8449}

% Ramka na regułę
\newtcolorbox{regula}[1][]{
  colback=regula, colframe=cel, boxrule=0.8pt, arc=2pt,
  left=6pt, right=6pt, top=4pt, bottom=4pt, title=#1,
  fonttitle=\bfseries, coltitle=white
}

% Ramka na pułapkę z kartkówki
\newtcolorbox{pulapka}[1][]{
  colback=pulapka, colframe=okres, boxrule=0.8pt, arc=2pt,
  left=6pt, right=6pt, top=4pt, bottom=4pt, title=#1,
  fonttitle=\bfseries, coltitle=white
}

% Nagłówek karty: numer, tytuł, cel
\newcommand{\naglowekKarty}[3]{%
  {\LARGE\bfseries\color{cel} Karta #1: #2}\\[2pt]
  {\large Cel: #3}\par
  \vspace{0.3em}\noindent\rule{\linewidth}{1pt}\par
}
 \begin{document}

% Silnik tectonic (XeTeX): polskie znaki przez fontspec, nie przez inputenc/T1
\usepackage{fontspec}
\setmainfont{Arial}
\usepackage[polish]{babel}
\usepackage[a4paper,margin=1.5cm]{geometry}
\usepackage{amsmath,amssymb}
\usepackage{tikz}
\usetikzlibrary{arrows.meta,positioning,calc,shapes.geometric}
\usepackage[most]{tcolorbox}
\usepackage{enumitem}
\usepackage{booktabs}
\usepackage{xcolor}

\pagestyle{empty}
\setlength{\parindent}{0pt}
\setlength{\parskip}{0.4em}

\definecolor{cel}{HTML}{1F4E79}
\definecolor{regula}{HTML}{E8F1FA}
\definecolor{pulapka}{HTML}{FDECEA}
\definecolor{okres}{HTML}{C0392B}
\definecolor{skonczony}{HTML}{1E8449}

% Ramka na regułę
\newtcolorbox{regula}[1][]{
  colback=regula, colframe=cel, boxrule=0.8pt, arc=2pt,
  left=6pt, right=6pt, top=4pt, bottom=4pt, title=#1,
  fonttitle=\bfseries, coltitle=white
}

% Ramka na pułapkę z kartkówki
\newtcolorbox{pulapka}[1][]{
  colback=pulapka, colframe=okres, boxrule=0.8pt, arc=2pt,
  left=6pt, right=6pt, top=4pt, bottom=4pt, title=#1,
  fonttitle=\bfseries, coltitle=white
}

% Nagłówek karty: numer, tytuł, cel
\newcommand{\naglowekKarty}[3]{%
  {\LARGE\bfseries\color{cel} Karta #1: #2}\\[2pt]
  {\large Cel: #3}\par
  \vspace{0.3em}\noindent\rule{\linewidth}{1pt}\par
}


\begin{document}

\naglowekKarty{5}{Zaokrąglanie (Mur i Policjant)}{zaokrąglać liczby dziesiętne i okresowe do wskazanego rzędu}

\vspace{0.3em}
\begin{regula}[REGUŁA – MUR I POLICJANT]
\textbf{Krok 1.} Postaw \textbf{mur} $|$ tuż za cyfrą, na której zaokrąglamy.\\
\textbf{Krok 2.} Spójrz na \textbf{pierwszą cyfrę za murem} (to jest policjant):
\begin{itemize}[leftmargin=*, itemsep=1pt]
  \item $0, 1, 2, 3, 4 \rightarrow$ \textbf{ŚPI}: cyfra przed murem zostaje bez zmian.
  \item $5, 6, 7, 8, 9 \rightarrow$ \textbf{ALARM}: dodaj $1$ do cyfry przed murem.
\end{itemize}
\textbf{Krok 3.} Odetnij wszystko za murem. Zera na końcu, które zostały po przecinku, zostają.
\end{regula}

\vspace{0.4em}
\textbf{Dlaczego tak? Druga droga: środek przedziału.} Liczby zaokrąglone do danego rzędu leżą co jeden krok na osi. Punkt środkowy między nimi to $5$ za murem. Wszystko poniżej środka zaokrąglamy w dół, wszystko od środka w górę.

\begin{center}
\begin{tikzpicture}[>=Stealth, x=10cm]
  % Przedział między 15,3 a 15,4 (zaokrąglanie do części dziesiątych)
  \draw[->] (-0.05,0) -- (1.05,0);
  \foreach \v/\lab in {0/{15,3}, 0.5/{15,35}, 1/{15,4}} {
    \draw (\v,0.1) -- (\v,-0.1);
    \node[below, font=\small] at (\v,-0.12) {$\lab$};
  }
  % Liczba 15,37 leży za środkiem: zaokrąglamy w górę
  \fill[okres] (0.7,0) circle (2.2pt) node[above=4pt, font=\small\bfseries, color=okres] {$15{,}37$};
  \draw[dashed, okres] (0.7,0) -- (0.7,-0.45);
  \node[font=\small, color=okres, anchor=north] at (0.7,-0.5) {$15{,}37 > 15{,}35 \Rightarrow 15{,}4$};
  \node[font=\small, anchor=west] at (-0.05,0.9) {Mur: $15{,}3\,|\,7$ $\;\Rightarrow\;$ policjant $7$ = ALARM $\Rightarrow 15{,}4$};
\end{tikzpicture}
\end{center}

\vspace{0.2em}
\textbf{Przykłady rozstrzygnięte:}
\begin{center}
\renewcommand{\arraystretch}{1.4}
\begin{tabular}{@{}lll@{}}
\toprule
Zadanie & Mur & Wynik \\
\midrule
$4{,}8\,|\,29$ do dziesiątych & policjant $2$ = ŚPI & $\mathbf{4{,}8}$ \\
$15{,}3\,|\,7$ do dziesiątych & policjant $7$ = ALARM & $\mathbf{15{,}4}$ \\
$0{,}35\,|\,1$ do setnych & policjant $1$ = ŚPI & $\mathbf{0{,}35}$ \\
$7{,}28\,|\,6$ do setnych & policjant $6$ = ALARM & $\mathbf{7{,}29}$ \\
$0{,}(6)=0{,}66\,|\,6\ldots$ do setnych & policjant $6$ = ALARM & $\mathbf{0{,}67}$ \\
$14{,}52$ do jedności: $14\,|\,{,}52$ & policjant $5$ = ALARM & $\mathbf{15}$ \\
$473{,}8$ do dziesiątek: $47\,|\,3{,}8$ & policjant $3$ = ŚPI & $\mathbf{470}$ \\
\bottomrule
\end{tabular}
\end{center}

\begin{pulapka}[PUŁAPKI Z KARTKÓWKI]
\begin{enumerate}[leftmargin=*, itemsep=2pt]
  \item \textbf{Okres najpierw rozpisz:} $0{,}(6)=0{,}666\ldots$. Dopiero wtedy stawiaj mur.
  \item \textbf{Efekt domina:} $3{,}998$ do setnych to $3{,}99\,|\,8$, policjant $8$ = ALARM, więc $3{,}99+0{,}01=\mathbf{4{,}00}$. Dwa zera po przecinku muszą zostać.
  \item \textbf{Patrz tylko na jedną cyfrę za murem.} Reszta nie ma znaczenia: $4{,}8\,|\,29$ to ŚPI, bo liczy się tylko $2$.
\end{enumerate}
\end{pulapka}

\end{document}
